% Shared preamble of the proof overview documents  <Folder>/<Name>_overview.tex.
% Each overview starts with
%     \input{../../scripts/overview_preamble.tex}
% (for ThinMatching/: % Shared preamble of the proof overview documents  <Folder>/<Name>_overview.tex.
% Each overview starts with
%     \input{../../scripts/overview_preamble.tex}
% (for ThinMatching/: % Shared preamble of the proof overview documents  <Folder>/<Name>_overview.tex.
% Each overview starts with
%     \input{../../scripts/overview_preamble.tex}
% (for ThinMatching/: \input{../scripts/overview_preamble.tex}) and is compiled from its own folder by
% scripts/build_overviews.py.
\documentclass[10pt,a4paper]{article}
\usepackage[T1]{fontenc}
\usepackage[utf8]{inputenc}
\usepackage[margin=2.2cm]{geometry}
\usepackage{amsmath,amssymb,amsthm}
\usepackage{array,booktabs,longtable}
\usepackage[dvipsnames]{xcolor}
\usepackage{enumitem}
\usepackage{tikz}
\usetikzlibrary{arrows.meta,positioning,shapes.multipart,fit,calc}
\usepackage[hidelinks]{hyperref}
\setlist{itemsep=2pt,topsep=4pt}
\renewcommand{\arraystretch}{1.2}

% Lean names, files and paper references
\newcommand{\lean}[1]{\textcolor{blue!50!black}{\texttt{\detokenize{#1}}}}
\newcommand{\file}[1]{\textcolor{OliveGreen!70!black}{\texttt{\detokenize{#1}}}}
\newcommand{\paperref}[1]{\textsf{#1}}
\newcommand{\hyp}[1]{\textcolor{BrickRed}{\textbf{#1}}}

% Boxes for the dependency diagrams
\tikzset{
  res/.style={draw, rounded corners, fill=blue!6, align=center, inner sep=4pt, font=\small},
  main/.style={res, fill=blue!18, very thick},
  cited/.style={res, fill=red!10, dashed},
  shared/.style={res, fill=gray!12},
  dep/.style={-{Stealth[length=2.2mm]}, thick},
}

\newcommand{\overviewtitle}[2]{%
  \begin{center}{\Large\bfseries #1}\\[4pt]{\small Proof overview of the Lean~4/Mathlib formalisation}\\[2pt]
  {\small #2}\end{center}}

% Standard paragraph describing the paper
\newcommand{\paperversion}{\noindent\textbf{Paper.} J.~Breitung, T.~Gervens, M.~Grohe.
  \emph{Graph Similarity Meets Optimal Transport}.\\ File \file{Graph_Similarity_Meets_Optimal_Transport.pdf} in the
  repository; statement, equation, page and line numbers refer to this PDF. \file{docs/NUMBERING.md} lists the Lean
  files and declarations of all numbered statements.\par\smallskip}
) and is compiled from its own folder by
% scripts/build_overviews.py.
\documentclass[10pt,a4paper]{article}
\usepackage[T1]{fontenc}
\usepackage[utf8]{inputenc}
\usepackage[margin=2.2cm]{geometry}
\usepackage{amsmath,amssymb,amsthm}
\usepackage{array,booktabs,longtable}
\usepackage[dvipsnames]{xcolor}
\usepackage{enumitem}
\usepackage{tikz}
\usetikzlibrary{arrows.meta,positioning,shapes.multipart,fit,calc}
\usepackage[hidelinks]{hyperref}
\setlist{itemsep=2pt,topsep=4pt}
\renewcommand{\arraystretch}{1.2}

% Lean names, files and paper references
\newcommand{\lean}[1]{\textcolor{blue!50!black}{\texttt{\detokenize{#1}}}}
\newcommand{\file}[1]{\textcolor{OliveGreen!70!black}{\texttt{\detokenize{#1}}}}
\newcommand{\paperref}[1]{\textsf{#1}}
\newcommand{\hyp}[1]{\textcolor{BrickRed}{\textbf{#1}}}

% Boxes for the dependency diagrams
\tikzset{
  res/.style={draw, rounded corners, fill=blue!6, align=center, inner sep=4pt, font=\small},
  main/.style={res, fill=blue!18, very thick},
  cited/.style={res, fill=red!10, dashed},
  shared/.style={res, fill=gray!12},
  dep/.style={-{Stealth[length=2.2mm]}, thick},
}

\newcommand{\overviewtitle}[2]{%
  \begin{center}{\Large\bfseries #1}\\[4pt]{\small Proof overview of the Lean~4/Mathlib formalisation}\\[2pt]
  {\small #2}\end{center}}

% Standard paragraph describing the paper
\newcommand{\paperversion}{\noindent\textbf{Paper.} J.~Breitung, T.~Gervens, M.~Grohe.
  \emph{Graph Similarity Meets Optimal Transport}.\\ File \file{Graph_Similarity_Meets_Optimal_Transport.pdf} in the
  repository; statement, equation, page and line numbers refer to this PDF. \file{docs/NUMBERING.md} lists the Lean
  files and declarations of all numbered statements.\par\smallskip}
) and is compiled from its own folder by
% scripts/build_overviews.py.
\documentclass[10pt,a4paper]{article}
\usepackage[T1]{fontenc}
\usepackage[utf8]{inputenc}
\usepackage[margin=2.2cm]{geometry}
\usepackage{amsmath,amssymb,amsthm}
\usepackage{array,booktabs,longtable}
\usepackage[dvipsnames]{xcolor}
\usepackage{enumitem}
\usepackage{tikz}
\usetikzlibrary{arrows.meta,positioning,shapes.multipart,fit,calc}
\usepackage[hidelinks]{hyperref}
\setlist{itemsep=2pt,topsep=4pt}
\renewcommand{\arraystretch}{1.2}

% Lean names, files and paper references
\newcommand{\lean}[1]{\textcolor{blue!50!black}{\texttt{\detokenize{#1}}}}
\newcommand{\file}[1]{\textcolor{OliveGreen!70!black}{\texttt{\detokenize{#1}}}}
\newcommand{\paperref}[1]{\textsf{#1}}
\newcommand{\hyp}[1]{\textcolor{BrickRed}{\textbf{#1}}}

% Boxes for the dependency diagrams
\tikzset{
  res/.style={draw, rounded corners, fill=blue!6, align=center, inner sep=4pt, font=\small},
  main/.style={res, fill=blue!18, very thick},
  cited/.style={res, fill=red!10, dashed},
  shared/.style={res, fill=gray!12},
  dep/.style={-{Stealth[length=2.2mm]}, thick},
}

\newcommand{\overviewtitle}[2]{%
  \begin{center}{\Large\bfseries #1}\\[4pt]{\small Proof overview of the Lean~4/Mathlib formalisation}\\[2pt]
  {\small #2}\end{center}}

% Standard paragraph describing the paper
\newcommand{\paperversion}{\noindent\textbf{Paper.} J.~Breitung, T.~Gervens, M.~Grohe.
  \emph{Graph Similarity Meets Optimal Transport}.\\ File \file{Graph_Similarity_Meets_Optimal_Transport.pdf} in the
  repository; statement, equation, page and line numbers refer to this PDF. \file{docs/NUMBERING.md} lists the Lean
  files and declarations of all numbered statements.\par\smallskip}

\begin{document}
\sloppy
\overviewtitle{Thin perfect matchings [HOG26]}{Library \file{ThinMatching/} \quad (Theorem 35 of the main paper)}

\noindent\textbf{Paper.} A.~Haqi, S.~Oveis Gharan. \emph{On Thin Perfect Matchings up to Polylogarithmic
Factors}, arXiv:2606.01330v2 (2026) [HOG26]. Statement, section, algorithm and equation numbers below refer
to [HOG26]; its numbering is independent of the main paper. The main paper (J.~Breitung, T.~Gervens,
M.~Grohe, \emph{Graph Similarity Meets Optimal Transport},
\file{Graph_Similarity_Meets_Optimal_Transport.pdf}) cites [HOG26, Theorem~1.4] as its
\textbf{Theorem~35} (p.~32, l.~1021--1022, Appendix~E.2.1) and uses it, through Theorem~38, for the
multiplicative upper bound of Theorem~11 (cut distance). This document is for a reviewer who wants to
check that Theorem~35 is proved in Lean, conditionally on the tree-cut sparsifier theorem of
R\"acke and Shah [RS14]. The line-by-line correspondence table is \file{ThinMatching/HOG26_correspondence.pdf}.
\par\smallskip

\noindent\textbf{Name clashes.} The Lean names follow [HOG26]: \lean{ThinMatching.Lemma36} is
[HOG26, Lemma~3.6] (not Lemma~36 of the main paper), \lean{ThinMatching.theorem13_support} is
[HOG26, Theorem~1.3] (not the main paper's Theorem~11, whose Lean name is \lean{GsimOT.theorem11_of_rackeShah}),
and \lean{ThinMatching.theorem14_bip} is [HOG26, Theorem~1.4].

\section{Statements}

\subsection{Theorem 35 of the main paper}
\emph{Thinness} (main paper, Eq.~(23), p.~32, l.~1014--1016): a permutation matrix $P\in\mathrm{Perm}(n)$ is
$\alpha$-thin with respect to $Q\in\mathrm{DS}(n)$ if
\[ \sum_{(i,j)\in(X\times\overline Y)\cup(\overline X\times Y)} P_{ij}\;\le\;\alpha
   \sum_{(i,j)\in(X\times\overline Y)\cup(\overline X\times Y)} Q_{ij}\qquad\text{for all }X,Y\subseteq[n]. \]
\emph{Theorem~35} ([HOG26, Theorem~1.4]): there is an absolute constant $C_{\mathrm{thin}}>0$ such that every
$Q\in\mathrm{DS}(n)$ admits an $\alpha$-thin permutation matrix with $\alpha\le C_{\mathrm{thin}}\ln(2n)\ln\ln(4n)$.

In Lean (main library, \file{GsimOT/Theorem11_CutUpperBound/Definitions.lean}) this is the proposition
\lean{GsimOT.ThinPermutationTheorem}, with thinness \lean{GsimOT.IsThin}: for $\pi:V\simeq W$ and all
$U_V\subseteq V$, $U_W\subseteq W$,
\[ \#\{v : \neg(v\in U_V\leftrightarrow\pi(v)\in U_W)\}\le a\sum_{v,w}[\neg(v\in U_V\leftrightarrow w\in U_W)]\,Q_{vw}, \]
which is Eq.~(23) with $X=U_V$, $Y=U_W$: the pair $(v,w)$ is counted iff $(v,w)\in(X\times\overline Y)\cup(\overline X\times Y)$.
\lean{ThinPermutationTheorem} asks for one $C_{\mathrm{thin}}>0$ such that for all finite types $V,W$ with
$|V|=|W|$ and all $Q$ with \lean{IsDS Q} there is $\pi$ with \lean{IsThin} $(C_{\mathrm{thin}}\ln(2|V|)\ln\ln(4|V|))\,Q\,\pi$.
This is exactly Theorem~35 (with $V,W$ of size $n$ instead of $[n]$; ``$\alpha\le$'' is equivalent to ``$\alpha=$''
because the right-hand side of Eq.~(23) is nonnegative). For $n=0$ the statement is trivial; for $n\ge1$ the
bound is positive, as required by ``Let $\alpha>0$''. So \lean{GsimOT.ThinPermutationTheorem} is Theorem~35.

\subsection{Formalised statements of [HOG26]}
Throughout, $x(\delta(S))=$ \lean{cutWeight x S}, $|M\cap\delta(S)|=$ \lean{crossCount M S}, and
\lean{IsThin} $\alpha\,x\,M$ means $|M\cap\delta(S)|\le\alpha\,x(\delta(S))$ for all $S\subseteq V$
(\file{Basic.lean}; [HOG26] quantifies over $\emptyset\ne S\subsetneq V$, which is equivalent since both
sides vanish for $S\in\{\emptyset,V\}$). $x\in\mathrm{PM}(G)$ is \lean{IsFracPM x}: a convex combination of
perfect matchings. In the bipartite case $x:L\to R\to\mathbb R$ with \lean{IsFracPMBip x} (nonnegative,
all row and column sums $1$, i.e.\ $x$ is doubly stochastic), $\pi:L\simeq R$, and \lean{IsThinBip}.

\begin{description}[leftmargin=0pt]
\item[Theorem 1.4, bipartite] \lean{ThinMatching.theorem14_bip}\,(hRS : \lean{RackeShahTheorem}): there is
  $C>0$ such that for all finite $L,R$ with $|L|=|R|=n$ and every $x$ with \lean{IsFracPMBip x} there is
  $\pi:L\simeq R$ (not necessarily inside $\mathrm{supp}(x)$) with
  $|M_\pi\cap\delta(S)|\le C\ln(2n)\ln\ln(4n)\,x(\delta(S))$ for all $S=S_L\cup S_R$.
  \textbf{This is the statement used for Theorem~35.}
\item[Theorem 1.4, general] \lean{ThinMatching.theorem14}\,(hRS): there is $C>0$ such that every $x\in\mathrm{PM}(G)$
  on $V$ has a $C\ln(2|V|)\ln\ln(4|V|)$-thin perfect matching $M$ of the complete graph on $V$
  ($=C\ln(4n)\ln\ln(8n)$ for $|V|=2n$).
\item[Theorem 1.3] \lean{ThinMatching.theorem13_support}\,(h36 : \lean{Lemma36}): there is $C>0$ such that every
  $x\in\mathrm{PM}(G)$ has a perfect matching $M\subseteq\mathrm{supp}(x)$ (\lean{InSupport}) that is
  $C|V|\ln(2|V|)$-thin.
\item[Lemma 1.2] \lean{ThinMatching.lemma12}\,(no hypotheses): for $n\ge1$ the graph of Fig.~1 on $2n+2$ vertices
  has $x\in\mathrm{PM}(G)$ such that every perfect matching $M\subseteq\mathrm{supp}(x)$ has a cut
  $\emptyset\ne S\subsetneq V$ with $x(\delta(S))>0$ and $|M\cap\delta(S)|\ge\frac n2\,x(\delta(S))$;
  corollary \lean{lemma12_not_thin}: no such $M$ is $\alpha$-thin for $\alpha<n/2$.
\end{description}
Lemma~1.2 and Theorem~1.3 are not needed for Theorem~35; they are formalised for completeness.

\subsection{The two explicit hypotheses}
\begin{itemize}
\item \hyp{\lean{RackeShahTheorem}} (\file{TreeSparsifier.lean}) = [HOG26, Theorem~2.2] = [RS14]: there is
  $C>0$ such that every weighted graph $x$ on $n\ge1$ vertices has a $C\ln(2n)\ln\ln(4n)$-tree-cut sparsifier
  (\lean{IsTreeCutSparsifier}, Definition~2.1: a weighted tree $T$ with Steiner nodes, $V\hookrightarrow\mathcal V$,
  and $x(\delta(S))\le\lambda_T(S)\le q\,x(\delta(S))$ for all $S$). [HOG26] cites this result without proof;
  its proof (R\"acke--Shah, building on R\"acke's hierarchical decompositions) is far outside the scope of
  this formalisation. $\ln(2n)\ln\ln(4n)$ replaces $\log n\log\log n$ so that the bound is positive for $n\ge1$.
\item \hyp{\lean{Lemma36}} (\file{SupportConstrained/Sampling.lean}) = [HOG26, Lemma~3.6], whose proof [HOG26]
  omits (``a simple Chernoff bound and a cut-counting argument''): there are $c,K>0$ such that for every
  weighted graph $y\le1$ on $W$, every $N\ge\max(1,|W|)$ and every integer $\alpha\ge cN\ln(2N)$: if all
  nonempty proper cuts have $y(\delta(S))\ge1/(4N)$, then with probability $\ge1-1/(4N)$ the sample $Y$
  ($\alpha$ independent copies of each edge $e$, each kept with probability $y_e$) satisfies
  $Y(\delta(S))\le K\alpha\,y(\delta(S))$ for all $S$. Only Theorem~1.3 uses it.
\end{itemize}
Both are Lean \texttt{def}s of type \texttt{Prop} that appear as arguments of the theorems, never as
\texttt{axiom}s: the dependence is visible in every statement, a false hypothesis cannot make other parts
of the library inconsistent, and \texttt{\#print axioms} shows only \texttt{propext}, \texttt{Classical.choice},
\texttt{Quot.sound}. \textbf{Theorem~35 depends only on \lean{RackeShahTheorem}}, not on \lean{Lemma36}.

\section{Where the proof lives}
\noindent Paths are relative to \file{ThinMatching/} unless stated otherwise.
\begin{longtable}{>{\raggedright\arraybackslash}p{5.8cm} >{\raggedright\arraybackslash}p{9.8cm}}
\toprule \textbf{File} & \textbf{Content} \\ \midrule\endhead \bottomrule\endfoot
\file{Basic.lean} & weighted graphs, cuts, perfect matchings (fixed-point-free involutions), thinness, $\mathrm{PM}(G)$; bipartite versions (Sections~1, 2) \\
\file{TreeSparsifier.lean} & rooted trees, subtrees, tree cuts, $\lambda_T$, Definition~2.1, hypothesis \lean{RackeShahTheorem} (Theorem~2.2) \\
\file{Load.lean} & loads and discrepancies (Definitions~4.1, 4.2, 4.4) \\
\file{LoadVsCut.lean} & Proposition~4.3 (\lean{prop43}) and the charging lemma \lean{exists_cut_edge} \\
\file{Discrepancy.lean} & Lemma~4.5 (\lean{lemma45_bip}, \lean{lemma45}) \\
\file{TreePairing.lean} & Algorithm~2, Lemma~4.6 (\lean{lemma46_bip}, \lean{lemma46}) \\
\file{Theorem14.lean} & Theorem~1.4 (\lean{theorem14_bip}, \lean{theorem14}); bipartite/general translation; constants \\
\file{LowerBound.lean} & Lemma~1.2, Fig.~1 \\
\file{SupportConstrained/}\newline\file{Prob.lean} & finite product spaces of Bernoulli coins (probability model of Algorithm~1) \\
\file{SupportConstrained/}\newline\file{Sampling.lean} & the sampling step, hypothesis \lean{Lemma36} \\
\file{SupportConstrained/}\newline\file{Trimming.lean} & Definition~3.1, Lemma~3.2 \\
\file{SupportConstrained/}\newline\file{TutteWitness.lean} & Tutte witnesses (Definition~3.3, Theorem~2.4 via Mathlib's \lean{SimpleGraph.tutte}), Lemma~3.4, Eq.~(1) \\
\file{SupportConstrained/}\newline\file{Lemma35.lean} & Lemma~3.5 \\
\file{SupportConstrained.lean} & Theorem~1.3 (\lean{theorem13_support}), Algorithm~1 \\
\file{GsimOT/Theorem11_CutUpperBound/}\newline\file{FromThinMatching.lean} & bridge (main library): Theorem~35 and Theorem~11 of the main paper assuming only [RS14] \\
\end{longtable}

\section{Proof layout}
\subsection{Theorem 1.4 and the bridge to the main paper}
\begin{center}
\begin{tikzpicture}[x=1cm,y=1cm, leaf/.style={res, font=\scriptsize}]
  \node[shared] (t11) at (0,0) {Main paper, Theorem 11\\\lean{GsimOT.theorem11_of_rackeShah}};
  \node[main] (t35) at (0,-1.5) {Main paper, Theorem 35\\\lean{GsimOT.thinPermutationTheorem_of_rackeShah}};
  \node[main] (t14b) at (-1.6,-3.2) {Theorem 1.4 (bipartite)\\\lean{theorem14_bip}};
  \node[res] (t14g) at (4.2,-3.2) {Theorem 1.4 (general)\\\lean{theorem14}};
  \node[leaf] (tr) at (-6.0,-3.2) {translation $L\oplus R$\\\lean{bipPM}, \lean{cutWeight_bipGraph},\\\lean{load_bipPM}, \lean{log_mul_loglog_le}};
  \node[cited] (rs) at (-6.0,-5.4) {Theorem 2.2 [RS14]\\\lean{RackeShahTheorem}};
  \node[res] (p43) at (-2.2,-5.4) {Proposition 4.3\\\lean{prop43}};
  \node[res] (l45) at (1.4,-5.4) {Lemma 4.5\\\lean{lemma45_bip},\\\lean{lemma45}};
  \node[res] (l46) at (5.2,-5.4) {Lemma 4.6 (Alg.~2)\\\lean{lemma46_bip},\\\lean{lemma46}};
  \node[leaf] (ce) at (-4.6,-7.6) {charging lemma\\\lean{exists_cut_edge}};
  \node[leaf] (ts) at (-1.5,-7.6) {single-edge tree cut\\\lean{treeCutWeight_subtree}};
  \node[leaf] (odd) at (2.1,-7.6) {odd sets\\\lean{one_le_crossCount_of_odd}};
  \node[leaf] (lam) at (6.0,-7.6) {laminarity\\\lean{subtree_laminar},\\\lean{subtree_ssubset_of_separates}};
  \draw[dep] (t11) -- (t35);
  \draw[dep] (t35) -- (t14b);
  \draw[dep] (t14b) -- (tr);
  \draw[dep] (t14b) -- (rs); \draw[dep] (t14b) -- (p43); \draw[dep] (t14b) -- (l45); \draw[dep] (t14b) -- (l46);
  \draw[dep] (t14g) -- (rs); \draw[dep] (t14g) -- (p43); \draw[dep] (t14g) -- (l45); \draw[dep] (t14g) -- (l46);
  \draw[dep] (p43) -- (ce); \draw[dep] (p43) -- (ts);
  \draw[dep] (l45) -- (odd); \draw[dep] (l46) -- (lam);
\end{tikzpicture}
\end{center}
\noindent Theorem~11 of the main paper also needs the rest of \file{GsimOT/Theorem11_CutUpperBound}
(Lemmas~36, 37, Theorems~38, 39, \dots{} of the main paper, \lean{GsimOT.theorem11_of_thin}); see the overview of that folder.

\subsection{Theorem 1.3 and Lemma 1.2}
\begin{center}
\begin{tikzpicture}[x=1cm,y=1cm, leaf/.style={res, font=\scriptsize}]
  \node[main] (t13) at (0,0) {Theorem 1.3 (Algorithm 1)\\\lean{theorem13_support}};
  \node[main] (l12) at (6.0,0) {Lemma 1.2 (Fig.~1)\\\lean{lemma12}};
  \node[res] (l32) at (-6.2,-2) {Lemma 3.2 (Def.~3.1)\\\lean{lemma32}};
  \node[res] (l35) at (-2.4,-2) {Lemma 3.5\\\lean{lemma35},\\\lean{witness_bound_le}};
  \node[res] (blk) at (1.6,-2) {union over blocks\\\lean{prob_bad_blocks_le}};
  \node[res] (good) at (5.6,-2) {thinness from blocks\\\lean{crossCount_le_of_good}};
  \node[leaf] (l34) at (-7.0,-4.4) {Lemma 3.4\\\lean{card_witnesses_le}};
  \node[leaf] (eq1) at (-4.0,-4.4) {Eq.~(1)\\\lean{one_le_witnessF_weight}};
  \node[shared, font=\scriptsize] (tut) at (-0.7,-4.4) {Theorem 2.4 (Tutte)\\\lean{exists_witness_of_no_pm}};
  \node[leaf] (prob) at (3.0,-4.4) {probability model\\\file{Prob.lean}, \file{Sampling.lean}};
  \node[cited] (l36) at (6.4,-4.4) {Lemma 3.6\\\lean{Lemma36}};
  \draw[dep] (t13) -- (l32); \draw[dep] (t13) -- (l35); \draw[dep] (t13) -- (blk); \draw[dep] (t13) -- (good);
  \draw[dep] (l35) -- (l34); \draw[dep] (l35) -- (eq1); \draw[dep] (l35) -- (tut); \draw[dep] (l35) -- (prob);
  \draw[dep] (blk) -- (l36); \draw[dep] (blk) -- (prob);
\end{tikzpicture}
\end{center}
\noindent Styles: thick = main results, dashed red = explicit hypotheses (cited results), grey = results
outside this part (Mathlib's Tutte theorem; the main paper's Theorem~11).

\section{The proof of Theorem 1.4 step by step}
The proof of [HOG26, Section~4] has three ingredients and a short assembly.

\paragraph{Step 0: trees and tree-cut sparsifiers (Section~2, Definition~2.1, Theorem~2.2).}
A weighted tree with Steiner nodes is a \lean{RootedTree N} (parent function, every vertex reaches the root),
edge weights $y:N\to\mathbb R$ (edge $\{v,\mathrm{parent}(v)\}$ has weight $y_v$) and an embedding
$e:V\hookrightarrow N$. \lean{treeCutWeight} is $y(\delta_T(\mathcal S))$, \lean{lambdaT} is
$\lambda_T(S)=\min\{y(\delta_T(\mathcal S)):\mathcal S\cap V\in\{S,V\setminus S\}\}$, and
\lean{IsTreeCutSparsifier} $q$ is Definition~2.1. Theorem~2.2 is the hypothesis \lean{RackeShahTheorem}.

\paragraph{Step 1: loads and discrepancies (Definitions~4.1, 4.2, 4.4; \file{Load.lean}).}
The tree path between $a$ and $b$ uses the edge $f_v=\{v,\mathrm{parent}(v)\}$ iff exactly one of $a,b$ lies in
\lean{subtree T v}. Hence $\mathrm{load}_M(f_v)$ (\lean{load}, \lean{loadBip}) counts the matching edges with
exactly one endpoint in the subtree of $v$, and the discrepancies are
$\Delta_{\mathrm{bip}}=\bigl||L\cap\cdot|-|R\cap\cdot|\bigr|$ (\lean{discBip}) and
$\Delta_{\mathrm{non\text{-}bip}}=\mathbf 1[|V\cap\cdot|\text{ odd}]$ (\lean{discNonBip}) of $V\cap\mathrm{subtree}(v)$.

\paragraph{Step 2: Proposition 4.3 (p.~7; \lean{prop43}).}
If $T$ is a $q$-tree-cut sparsifier of $x$ and $\mathrm{load}_M(f_v)\le\beta\,x(\delta(V\cap\mathrm{subtree}(v)))$
for every tree edge, then $M$ is $\beta q$-thin. Lean follows the paper: take an optimal tree cut $\mathcal S'$
for $S$ (\lean{exists_eq_ciInf_of_finite}); every crossing matching edge is charged to an edge
$f_v\in\delta_T(\mathcal S')$ on its tree path (\lean{exists_cut_edge}); since the single-edge cut
$\mathrm{subtree}(v)$ has value $y_v$ (\lean{treeCutWeight_subtree}) and the sparsifier does not underestimate
cuts, $\mathrm{load}_M(f_v)\le\beta x(\delta(\cdot))\le\beta\lambda_T(\cdot)\le\beta y_v$; summing,
$|M\cap\delta(S)|\le\beta\,y(\delta_T(\mathcal S'))=\beta\lambda_T(S)\le\beta q\,x(\delta(S))$.

\paragraph{Step 3: Lemma 4.5 (p.~8; \file{Discrepancy.lean}).}
Bipartite (\lean{lemma45_bip}): by the degree constraints,
$\bigl||S_L|-|S_R|\bigr|=|x(S_L,R)-x(L,S_R)|\le x(\delta(S))$, as in the paper.
General (\lean{lemma45}): $x(\delta(S))\ge1$ for $|S|$ odd; Lean writes $x$ as a convex combination of
perfect matchings, each of which crosses an odd set (\lean{one_le_crossCount_of_odd},
\lean{cutWeight_eq_sum_crossCount}), instead of invoking the odd-set constraint (Theorem~2.5).

\paragraph{Step 4: Lemma 4.6, Algorithm 2 (pp.~9--10; \file{TreePairing.lean}).}
For points placed at tree vertices there is a perfect matching with $\mathrm{load}_M(f)=\Delta(f)$ for every
tree edge (\lean{lemma46_bip} for $|L|=|R|$, \lean{lemma46} for $|V|$ even). Lean runs the bottom-up greedy
one pair at a time: it matches a matchable pair inside a subtree of minimal size that contains one (the
first pair the greedy forms) and recurses (strong induction on the number of points). The paper's key
observation (inside the chosen subtree no proper subtree contains a matchable pair) becomes
\lean{subtree_laminar}, \lean{subtree_ssubset_of_separates}: every edge separating the pair sees load and
discrepancy both grow by one, all other edges see no change.

\paragraph{Step 5: assembly (Section~4; \file{Theorem14.lean}).}
\lean{theorem14}: obtain the sparsifier $T$ of $x$ from \lean{RackeShahTheorem} with $q=C\ln(2|V|)\ln\ln(4|V|)$,
the matching $M$ of \lean{lemma46} ($|V|$ is even because a perfect matching exists), check the hypothesis of
\lean{prop43} with $\beta=1$ by \lean{lemma45}, and conclude that $M$ is $q$-thin.
\lean{theorem14_bip}: the same argument on the weighted graph \lean{bipGraph x} on $L\oplus R$, with the
translations \lean{bipPM} ($\pi\mapsto$ involution of $L\oplus R$), \lean{crossCount_bipPM},
\lean{cutWeight_bipGraph}, \lean{load_bipPM}. Since $|L\oplus R|=2n$ the sparsifier bound is
$C\ln(4n)\ln\ln(8n)\le6C\ln(2n)\ln\ln(4n)$ (\lean{log_mul_loglog_le}), so \lean{theorem14_bip} holds with
constant $6C$, where $C$ is the constant of \lean{RackeShahTheorem}. The case $n=0$ is trivial.

\section{The proof of Theorem 1.3 (Section 3) in brief}
\lean{theorem13_support} follows Algorithm~1 with $n=|V|$, $\varepsilon=1/(4n)$ and
$\alpha=\lceil c'n\ln(2n)\rceil$, $c'=\max(c,4)$: (1) trimming (Definition~3.1, \lean{lemma32}): blocks that are
$\varepsilon$-edge-connected, removed weight $\le(n-1)\varepsilon\le1/4$; (2) sampling (\lean{sampleProb});
(3) Lemma~3.5 (\lean{lemma35}, via Tutte's theorem \lean{exists_witness_of_no_pm}, at most $(n+1)^n$ witnesses
\lean{card_witnesses_le}, and $x(F)\ge1$ \lean{one_le_witnessF_weight}): no perfect matching with probability
$\le1/4$ (\lean{witness_bound_le}); (4) \lean{Lemma36} on each block plus a union bound over $\le n$ blocks
(\lean{prob_bad_blocks_le}): some block bad with probability $\le1/4$; (5) an outcome of positive probability
avoiding both events (\lean{exists_not_of_prob_lt_one}) gives $M\subseteq Y\subseteq\mathrm{supp}(x)$, and
\lean{crossCount_le_of_good} (with \lean{sum_cut_eq_sum_blocks}) gives thinness for every cut, also cuts through
several blocks. Constant: $C=K(\max(c,4)+2)$.

\section{The bridge to the main paper}
\file{GsimOT/Theorem11_CutUpperBound/FromThinMatching.lean} (main library; it imports
\file{ThinMatching.Theorem14}):
\begin{itemize}
\item \lean{GsimOT.thinPermutationTheorem_of_rackeShah} (hRS : \lean{ThinMatching.RackeShahTheorem}) :
  \lean{GsimOT.ThinPermutationTheorem}. Proof: take $C$ from \lean{theorem14_bip}; for $Q$ with \lean{IsDS Q}
  the triple (\lean{nonneg}, \lean{row_sum}, \lean{col_sum}) is \lean{IsFracPMBip Q}; the resulting $\pi$ is thin
  in the sense of \lean{GsimOT.IsThin} after unfolding \lean{bipCrossCount}, \lean{bipCutWeight} (a \texttt{simpa}).
  Hence Theorem~35 holds with $C_{\mathrm{thin}}=6\,C_{\mathrm{RS}}$, where $C_{\mathrm{RS}}$ is the constant of
  \lean{RackeShahTheorem}.
\item \lean{GsimOT.theorem11_of_rackeShah}, \lean{GsimOT.theorem11_power_of_rackeShah}: Theorem~11 of the main
  paper (both bounds) assuming only [RS14], by plugging the above into
  \lean{GsimOT.theorem11_of_thin} and \lean{GsimOT.theorem11_power_of_thin}.
\end{itemize}

\section{Differences from [HOG26]}
\begin{itemize}
\item \emph{Representation.} Trees are rooted (parent function); this is no loss of generality. Loads and
  discrepancies are taken on the subtree side $V_s(f)$; the paper uses the root side $V_r(f)$, which has the same
  discrepancy since $|L|=|R|$ resp.\ $|V|$ is even. Perfect matchings are fixed-point-free involutions (general)
  or bijections $L\simeq R$ (bipartite). $\mathrm{PM}(G)$ is defined as the convex hull of perfect matchings
  (bipartite: degree constraints, equivalent by Birkhoff--von Neumann).
\item \emph{Constants.} $\log n\log\log n$ is written $\ln(2n)\ln\ln(4n)$ (positive for $n\ge1$), as in the main
  paper's Theorem~35. The bipartite constant is $6C_{\mathrm{RS}}$ because the sparsifier lives on $2n$ vertices.
\item \emph{Lemma 4.5 (general)} uses only the easy direction of Edmonds' theorem (Theorem~2.5).
\item \emph{Lemma 4.6} is proved by matching one pair at a time in a minimal subtree, not by running the greedy
  vertex by vertex; the invariant is the paper's.
\item \emph{Theorem 1.3}: $\varepsilon=1/(4|V|)$ (the paper's $1/(8n)$ for $|V|=2n$); $x(F)\ge1$ instead of
  $x(F)\ge1/2$; trimming by block labellings (one new block per step); the paper argues thinness only per
  component, Lean also for cuts through several components (no loss). \lean{Lemma36} is formulated with a global
  $N\ge|W|$ and failure probability $\le1/(4N)$, which is what the omitted proof gives and what the proof of
  Theorem~1.3 uses.
\item \emph{Not formalised}: Corollary~1.6 (metric distortion; needs [ACR23]) and Theorem~2.3 (algorithmic
  sparsifiers, not used).
\end{itemize}
No error in [HOG26] was found in the formalised parts.

\section{How to check}
From the repository root (with \texttt{elan} on the \texttt{PATH}):
\begin{itemize}
\item \texttt{lake build ThinMatching} builds the library (it imports only Mathlib);
  \texttt{lake build GsimOT.Theorem11\_CutUpperBound.FromThinMatching} builds the bridge.
\item \texttt{bash build.sh} builds everything, greps for \texttt{sorry}/\texttt{admit}/\texttt{axiom}, and runs
  \file{scripts/CheckAxioms.lean}, which contains
  \texttt{\#print axioms} for \lean{ThinMatching.theorem14_bip}, \lean{ThinMatching.theorem14},
  \lean{ThinMatching.theorem13_support}, \lean{ThinMatching.lemma12}, the lemmas of Sections~3--4 and
  \lean{GsimOT.thinPermutationTheorem_of_rackeShah}, \lean{GsimOT.theorem11_of_rackeShah}. Each must report
  only \texttt{propext}, \texttt{Classical.choice}, \texttt{Quot.sound}.
\item To review the statements: read the definitions in \file{Basic.lean} and \file{TreeSparsifier.lean}, the
  hypothesis \lean{Lemma36} in \file{SupportConstrained/Sampling.lean}, the main theorems in
  \file{Theorem14.lean}, \file{SupportConstrained.lean}, \file{LowerBound.lean}, and the bridge with
  \lean{GsimOT.IsThin}, \lean{GsimOT.ThinPermutationTheorem} in \file{GsimOT/Theorem11_CutUpperBound/Definitions.lean}.
\end{itemize}
\end{document}
